\par\medskip
\begin{center}
\textls[90]{b) Das Jahr 1808. Beginn der Korrespondenz mit \textsc{Schumacher}.}
\end{center}
\noindent\hspace*{\parindent}Der Briefwechsel zwischen \textsc{Gauss} und \textsc{Schumacher}, der sich von 1808—1850,
also über einen Zeitraum von 42 Jahren erstreckt, wird durch einen
Brief \textsc{Schumachers} vom 2. April 1808 eingeleitet, in dem \textsc{Schumacher} \textsc{Gauss}
auffordert, die Aufgabe zu lösen, die der Spanier \textsc{Pedrayes} gestellt hatte. Wir
wollen zunächst über diese Aufgabe und über ihren Autor einige Angaben
machen, natürlich nur soweit sie zum Verständnis des folgenden erforderlich
sind.
\par
\textsc{Augustinus Pedrayes} (1744—1817) war 1769—1791 Lehrer am Pagenkollegium
(später adliges Seminar genannt) in Madrid; er veröffentlichte 1796
eine Aufgabe in der Form eines Programms\footnote{Erschienen in lateinischer und spanischer Sprache: Ex typographia Regia in Madrid; abgedruckt lateinisch im 9. Hefte 1799 von \textsc{Hindenburgs} Archiv der reinen und angewandten Mathematik, S. 85 ff. und in deutscher Übersetzung von H. \textsc{Murhardt} in den Göttinger Nachrichten von 1798.}, die im wesentlichen folgendermassen
lautet:\par
\clearpage\SourceLeaf{211}\ArithmeticHead{126}{L}{L. SCHLESINGER, ÜBER GAUSS' FUNKTIONENTHEORETISCHE ARBEITEN.}
\indent »\textls[100]{Problem}: die Integralgleichung zu finden, die dem folgenden Differential entspricht:
\[
\begin{aligned}
&\frac{ar^2\,dx}{\sqrt{(r-x)x}}+\frac{br^2\,dx}{\sqrt{(4r-x)x}}
+\frac{cr^2\,du}{\sqrt{(r-u)u}}+\frac{er^2\,du}{\sqrt{(4r-u)u}}\\[5pt]
&+\frac{fr\,dx\sqrt{4r^2-rx}}{\sqrt{(r-x)x}}
+\frac{hr\,dx\sqrt{r^2-rx}}{\sqrt{(4r-x)x}}
+\frac{kr\,du\sqrt{4r^2-ru}}{\sqrt{(r-u)u}}\\[5pt]
&+\frac{gr\,du\sqrt{r^2-ru}}{\sqrt{(4r-u)u}}
+\frac{lr\,dx\sqrt{4r-u}}{\sqrt{x}}
+\frac{mr\,du\sqrt{4r-x}}{\sqrt{u}}
+\frac{nr\,dx\sqrt{r-u}}{\sqrt{x}}\\[5pt]
&+\frac{pr\,du\sqrt{r-x}}{\sqrt{u}}
+\frac{qr\,dx\sqrt{(4r-u)(r-u)}}{\sqrt{rx}}
+\frac{sr\,du\sqrt{(4r-x)(r-x)}}{\sqrt{ru}}\\[5pt]
&+\frac{tru\,dx}{\sqrt{rx}}+\frac{zrx\,du}{\sqrt{ru}}=dY.\text{«}
\end{aligned}
\]
\noindent Hierin bedeuten \(r,a,b,c,\ldots,z\) Konstanten.
\par
Dem Abdruck in \textsc{Hindenburgs} Archiv sind Bemerkungen von \textsc{Wilhelm
Pfaff} (Professor der Mathematik in Dorpat, später in Erlangen) angeschlossen,
die eine Lösung der \textsc{Pedrayesschen} Aufgabe von J. Fr. \textsc{Pfaff} (dem uns bereits
bekannten, \textsc{Gauss} befreundeten Helmstedter Professor, \textsc{Wilhelms} Bruder)
enthalten. Es heisst daselbst, dass »die Integration des 1., 2., 3., 4. Teiles
von der Quadratur des Kreises, des 5., 6. von der Rektifikation der Ellipse,
des 7., 8. von der Rektifikation der Hyperbel abhängt\footnote{\textsc{Pedrayes} hebt dies in seinen Erläuterungen zu dem Problem selbst hervor.}; die übrigen Teile
können angesehen werden als Differentiale von Flächen krummer Linien, deren
Beschaffenheit sich nach dem Verhalten zwischen \(u\) und \(x\) richtet. Es kommt
nun eigentlich — wie ich glaube — bei der Aufgabe darauf an, \(u\) und \(x\)
durch einander oder durch eine dritte Grösse \(\varphi\) so zu bestimmen, dass, obgleich
die eigentlichen Integrale transzendent sind, doch das gesamte Integral
algebraisch werde. Dies wird nicht bei allen Werten \(a,b,\ldots,z\) tunlich sein,
daher müssen noch ferner die Werte dieser Konstanten bestimmt werden, bei
welchen es angeht.« — Die Lösung von J. F. \textsc{Pfaff} bezieht sich auf \(u=x\).
\par
Der Pariser Akademie wurden eingereicht eine aus Berlin stammende
Lösung, die aber als unzureichend erklärt wurde, und die Lösung von \textsc{Pedrayes}
selbst, in bezug auf die beschlossen wurde, das Urteil der Akademie
nicht zu veröffentlichen. 1805 erschien dann die \textsc{Pedrayessche} Lösung unter
dem Titel: »Opusculum primum. Solutio problematis propositi anno 1797
[sic!]; a subscriptorum societate literaria pervulgata.«
\par
In seinem Briefe vom 2. April 1808 (Werke X, 1, S. 242) erwähnt \textsc{Schumacher}
sowohl die Auflösung von \textsc{Pfaff} als auch die eigene Auflösung des
\textsc{Pedrayes}. Auch \textsc{Gauss} bezieht sich in seinem Antwortschreiben vom 17. Sept.\par
\clearpage\SourceLeaf{212}\ArithmeticHead{127}{R}{VIb. KORRESPONDENZ MIT SCHUMACHER 1808.}
\noindent 1808 (Werke X, 1, S. 243) auf die von \textsc{Pfaff} gegebene Lösung und fährt
dann fort:
\par
»Vielleicht wäre ich im Besitz von Wahrheiten, die zur Entscheidung
dieser Sache dienen könnten \mbox{.\ .\ .\ .}\quad Mit Kreisfunktionen und Logarithmischen
wissen wir jetzt umzugehen, wie mit dem 1 mal 1, aber die herrliche Goldgrube,
die das Innere der höheren Funktionen enthält, ist noch fast ganz
Terra Incognita. Ich habe darüber ehemals sehr viel gearbeitet und werde
dereinst ein eigenes grosses Werk darüber geben, wovon ich bereits in meinen
\textit{Disquisitiones arithm.} p. 593\footnote{Werke I, S. 412, 413; Art. 335. Es ist vielleicht angebracht, diese berühmte Stelle hier wörtlich wiederzugeben. Sie lautet:\par
»Im übrigen reichen die Prinzipien der Theorie der Kreisteilungsgleichungen, die wir im Begriffe sind darzulegen, sehr viel weiter, als hier entwickelt werden kann. Sie können nicht nur auf die Kreisfunktionen, sondern mit dem gleichen Erfolge auf viele andere transzendente Funktionen angewandt werden, z. B. auf solche, die von dem Integral
\[
\int\frac{dx}{\sqrt{1-x^4}}
\]
abhängen\mbox{.\ .\ .\ .\ .}\quad Da wir aber über diese transzendenten Funktionen ein umfangreiches Werk vorbereiten, \mbox{.\ .\ .\ .\ .}\ so schien es angebracht, hier nur die Kreisfunktionen allein zu betrachten.«}
einen Wink gegeben habe. Man gerät in Erstaunen
über den überschwenglichen Reichtum an neuen höchst interessanten
Wahrheiten und Relationen, die dergleichen Funktionen darbieten (wohin u. a.
auch diejenigen gehören, mit denen die Rektifikation der Ellipse und Hyperbel
zusammenhängt). Es könnte wohl sein, dass gerade aus diesen Untersuchungen
die Beantwortung der \textsc{Pedrayes}-Aufgabe sich entnehmen liesse, vorausgesetzt,
dass sie eine Auflösung zulässt, die wirklich einen Wert hat: allein wenn ich
auch klarer sähe, dass die ganze Aufgabe zu etwas führen könnte, als das
bis jetzt der Fall ist, würde ich doch jetzt von dieser Untersuchung abstrahieren
müssen, da ich mich erst dann in diese weitaussehende Materie wieder
hinein werfen werde, wenn ich an die Ausarbeitung jenes grossen Werkes
werde denken können. Dazu bin ich aber jetzt noch mit zu vielen andern,
mir nicht minder interessanten Untersuchungen überhäuft.«
\par
Es ist wohl anzunehmen, dass trotz der Schlusswendung, die \textsc{Gauss} in
dieser brieflichen Äusserung gebraucht, die Aufforderung \textsc{Schumachers}, die
Lösung der \textsc{Pedrayesschen} Aufgabe zu versuchen, den äussern Anstoss dazu
gegeben hat, dass \textsc{Gauss} 1808 die Beschäftigung mit den elliptischen Transzendenten\par
\clearpage\SourceLeaf{213}\ArithmeticHead{128}{L}{L. SCHLESINGER, ÜBER GAUSS' FUNKTIONENTHEORETISCHE ARBEITEN.}
\noindent wieder aufnahm; jedenfalls bringen die Jahre 1808 und 1809 wieder
eine Reihe wichtiger Beiträge zu der in Rede stehenden Theorie.
\par
Im Handbuch 16, Bb, beginnt auf Seite 40 eine umfangreiche, in deutscher
Sprache redigierte Abhandlung mit der Überschrift: \textit{Zur Theorie der
transzendenten Funktionen gehörig}; die Seiten 40—51 des Handbuchs sind Werke
X, 1, S. 287—307 (Artikel 1—16), die Seiten 51—53 des Handbuchs, Werke
III, S. 442—445 (Artikel 8, 9) abgedruckt\footnote{Werke III, S. 436—442 (Artikel 1—7) hat \textsc{Schering} nur einen Auszug der Seiten 40—51 des Handbuchs gegeben; es war darum nicht zu umgehen, den vollständigen Text dieser Seiten im Bande X, 1 der Werke abzudrucken; die Werke III, S. 442—445 abgedruckten Artikel 8, 9, geben den Text der Handschrift mit so unerheblichen Auslassungen, dass von einem Wiederabdruck abgesehen werden konnte.}. Die Abfassungszeit dieser Abhandlung
ist, wie auch \textsc{Schering}, Werke III, S. 494 bemerkt\footnote{Vergl. auch die Bemerkung Werke X, 1, S. 320—323.}, durch eine
für \textsc{Jacobi} bestimmte Mitteilung von \textsc{Gauss} an \textsc{Schumacher}\footnote{Vgl. \textsc{Jacobis} Werke I, S. 394, Brief an \textsc{Legendre}, wo es heisst: M. \textsc{Gauss} ayant appris de celui-ci [\textsc{Jacobis} Arbeiten in den Astronomischen Nachrichten von 1827, \textsc{Jacobis} Werke I, S. 29—48] m’a fait dire qu’il avait développé déjà en 1808 les cas de 3 sections, 5 sections, et de 7 sections, et trouvé en même temps les nouvelles échelles de modules qui s’y rapportent. Cette nouvelle, à ce qui me paraît, est bien intéressante.« Vergl. \textsc{Legendre} ebenda S. 398, \textsc{Jacobi} ebenda S. 416, \textsc{Legendre} ebenda S. 418, 428.} vom Juni 1827
auf 1808 festgelegt. Sie beginnt mit der Entwicklung der schon oben (S. 96,
Gl. (15)) erwähnten Formel für die \textls[100]{lineare Transformation der von
beiden Argumenten abhängigen Thetafunktion} \(\displaystyle T=\sum_{k=-\infty}^{+\infty}e^{-\alpha(k+\omega)^2}\),
oder, mit Benutzung unserer obigen Bezeichnung \(\vartheta_{00}(\psi|x)\),
\[
T=e^{-\alpha\omega^2}\vartheta_{00}\left(\frac{\alpha\omega}{\pi i}\biggm|e^{-\alpha}\right),
\]
die \textsc{Gauss} erhält, indem er diese Funktion \(T\) in eine nach Kosinus der Vielfachen
von \(2\pi\omega\) fortschreitende Reihe entwickelt und die Koeffizienten in
der heute nach \textsc{Fourier} benannten Weise durch Integration zwischen den
Grenzen \(\omega=0\) bis \(\omega=1\) bestimmt. \textsc{Gauss} bezeichnet im Art. [8] (Werke
III, S. 442) die Theorie der linearen Transformation als »die schönen Lehrsätze
der Reziprozität«, wohl weil es sich dabei um den Übergang von
\[
x=e^{-\frac{m}{n}\pi},\ y=e^{\frac{\lambda}{n}\pi}
\]
zu
\[
x'=e^{-\frac{n}{m}\pi},\ y'=e^{\frac{\lambda}{m}\pi i}
\]
\clearpage\SourceLeaf{214}\ArithmeticHead{129}{R}{VIc. DER BILINEARE ALGORITHMUS.}
\noindent handelt. Die Artikel 2—14 (Werke X, 1, S. 290—306) beziehen sich auf
die Umformungen der Produkte in Reihen, wobei die schon in der Scheda
Ac aufgestellte fundamentale Identität aufs neue abgeleitet wird. Da \textsc{Gauss}
diese Umformungen in einer späteren Aufzeichnung (\textit{Hundert Theoreme etc.},
Werke III, S. 461) ausführlich behandelt, gehen wir an dieser Stelle nicht
näher darauf ein. Der Art. 15 (ebenda, S. 306) behandelt die Siebenteilung,
der Art. 16 (ebenda, S. 307) die Dreiteilung der Perioden, was mit der oben
(Fussnote \textsuperscript{184)}) wiedergegebenen Mitteilung, die \textsc{Gauss} an \textsc{Jacobi} hatte gelangen
lassen, übereinstimmt\footnote{Auf den Seiten 63, 72, 73 des Handbuchs sind die Formeln über lemniskatische Funktionen aufgezeichnet, die Werke III, S. 405, 406, 409—412 untermischt mit früheren Notizen wiedergegeben sind, und über die wir oben S. 123 berichtet haben.}. Die Fünfteilung hatte \textsc{Gauss}, wie der Art. 14
(ebenda, S. 303—306) zeigt, 1808 zwar in Angriff genommen, aber, wie wir
sehen werden, erst im folgenden Jahre ausgeführt.
